\chapter{Yöntem ve Sistem Tasarımı}
\label{ch:yontem}

\section{Sistem Mimarisi}
\label{sec:mimari}

Sistem, fiziksel ve mantıksal olarak üç bölgeden oluşur:

\begin{enumerate}[label=\textbf{\Alph*.},itemsep=2pt]
    \item \textbf{Araç İçi Modül}: ESP32-CAM (kamera + Wi-Fi),
          STM32 NUCLEO-F767ZI (olay yöneticisi), HM-10 BLE köprü modülü,
          harici LED'ler, aktif buzzer, TARA ve PANİK butonları.
    \item \textbf{Bulut Tanıma}: AWS EC2 üzerinde Docker konteynerinde
          çalışan Flask + InsightFace hizmeti, pickle veri tabanı.
    \item \textbf{Acil Çağrı}: Sürücünün BLE üzerinden eşli Android
          telefonu, üzerinde çalışan TaxiGuard uygulaması.
\end{enumerate}

\Cref{fig:sistem-mimarisi} bu üç bölge arasındaki veri akışını gösterir.

\begin{figure}[h!]
\centering
\begin{tikzpicture}[
    node distance=10mm and 10mm,
    box/.style={draw=ink, thick, minimum width=28mm, minimum height=12mm,
                align=center, font=\small\sffamily},
    cloud/.style={draw=ink, thick, minimum width=34mm, minimum height=12mm,
                  align=center, font=\small\sffamily, fill=ink!5},
    arr/.style={-{Latex[length=2mm]}, thick, ink},
    arracc/.style={-{Latex[length=2mm]}, thick, accent},
    arrsig/.style={-{Latex[length=2mm]}, thick, signal},
    lab/.style={font=\scriptsize\ttfamily, ink},
]
\node[box] (esp) {ESP32-CAM\\\scriptsize 10-frame burst, Wi-Fi};
\node[box, below=of esp] (stm) {STM32 F767ZI\\\scriptsize Olay yöneticisi};
\node[box, right=22mm of esp] (wifi) {Wi-Fi 2.4\,GHz};
\node[cloud, right=of wifi] (ec2) {AWS EC2\\\scriptsize Flask + InsightFace};
\node[box, below=of ec2] (phone) {Android\\\scriptsize Intent.ACTION\_CALL};
\node[box, left=of phone] (hm10) {HM-10 BLE};

\draw[arr] (esp.east) -- node[lab, above]{HTTP POST · JPEG} (wifi.west);
\draw[arr] (wifi.east) -- (ec2.west);
\draw[arracc] (ec2.west) to[bend left=30] node[lab, below]{JSON · match} (wifi.east);
\draw[arr] (esp) -- node[lab, right]{UART1\\\scriptsize CAPTURE/RESULT} (stm);
\draw[arr] (stm.east) -- node[lab, above]{UART2} (hm10.west);
\draw[arr] (hm10.east) -- node[lab, above]{BLE notify} (phone.west);
\draw[arrsig] (phone.south) -- ++(0,-7mm) node[lab, right]{GSM · 155};
\end{tikzpicture}
\caption{Sistem üst seviye blok şeması: araç içi modül, bulut tanıma ve acil çağrı katmanları.}
\label{fig:sistem-mimarisi}
\end{figure}

Mimari tercihinde üç temel ilke belirleyici olmuştur:

\begin{itemize}[itemsep=2pt]
    \item \textbf{Tek sorumluluk:} Her bileşen, yalnızca tek bir
          görevden sorumludur (ESP32-CAM: kamera+ağ, STM32: olay yönetimi,
          EC2: tanıma, telefon: çağrı). Bu ayrışma, bir bileşenin
          değişmesinin diğerlerini etkilememesini sağlar.
    \item \textbf{Köprü olarak gömülü:} STM32, kamera ve internet
          erişiminin doğrudan tarafı değildir; ESP32-CAM bu görevi
          devralır. Bu, kamera/Wi-Fi'nin başarısız olması durumunda STM
          tarafındaki olay zincirinin (panik, BLE) etkilenmemesini garanti
          eder.
    \item \textbf{Bulut tarafı içe dönük:} Tanıma sunucusu yalnızca
          tek bir endpoint (\texttt{POST /search}) sunar; tüm karmaşık
          mantık (kalite filtresi, multi-frame agregasyon) burada
          gizlenir, gömülü taraf yalnızca ham JPEG verisi gönderir.
\end{itemize}

\section{Tanıma Yaklaşımı: Alternatiflerin Değerlendirilmesi}
\label{sec:alt-degerlendirme}

Gelişme raporunda \cite{ekin_gelisme_2026} başlatılan ``API destekli mi,
yerel mi?'' tartışması, kapanışında üç alternatifi karşılaştırarak
gerçekleştirme için API destekli yaklaşımı seçmiştir
(\Cref{tab:alternatif-degerlendirme}). Bu seçimin önemli bir nüansı,
``API destekli''nin \emph{kapalı kutu üçüncü taraf bulut servisi}
(AWS Rekognition, Azure Face, Face++) değil; ekibin kendi sunucusuna
deploy ettiği \emph{açık kaynak InsightFace REST API} olarak
yorumlanmasıdır. Bu yorum sayesinde:

\begin{itemize}[itemsep=2pt]
    \item Kullanılan model (\texttt{buffalo\_l}) ve eşik değer
          ($\tau$) ekibin kontrolündedir; tezde FAR/FRR ölçümleri
          dolgun bir savunma için ölçülmüştür.
    \item KVKK uyumu, veri akışını ekibin sahip olduğu sunucuda
          tuttuğu için pratik olarak sağlanmıştır.
    \item Kapatılma riski ya da fiyat değişikliği gibi üçüncü taraf
          bağımlılıkları yoktur.
\end{itemize}

\begin{table}[h!]
\centering
\caption{Yüz tanıma yaklaşımı için alternatif değerlendirme matrisi.}
\label{tab:alternatif-degerlendirme}
\small
\begin{tabularx}{\linewidth}{l X X l}
\toprule
\textbf{Alternatif} & \textbf{Avantaj} & \textbf{Dezavantaj/Risk} & \textbf{Karar}\\
\midrule
MCU üzerinde tanıma (Pico/F767) & Düşük maliyet, internet bağımsız &
Hesap gücü/bellek yetmez; doğruluk düşer & Elendi \\
Yerel ağır model (Pi 4 + CNN) & İnternet bağımsız, gecikme yerel &
Donanım maliyeti, ısı, takvim riski & Ertelendi \\
Kapalı üçüncü-taraf API (AWS Rekognition) & Hızlı entegrasyon &
KVKK belirsizliği, fiyat/kapatma riski, model üzerinde kontrol yok & Elendi \\
\textbf{Kendi InsightFace sunucumuz} & Tam kontrol, KVKK uyumlu, ücretsiz &
Sunucu işletim yükü (EC2 \euro 5/ay) & \textbf{Seçildi} \\
\bottomrule
\end{tabularx}
\end{table}

\section{Çoklu-Kare Burst Akışı}

Tek bir karelik snapshot ile tanıma, geçici tıkanma, hareket bulanıklığı
veya kırpılan yüz açısı durumlarında kırılgan davranır.  Bu projede,
ESP32-CAM tetikleme komutu alır almaz \textbf{5 FPS hızında 2 saniye
boyunca 10 kare} yakalar ve her birini ardışık olarak sunucuya iletir
(\Cref{fig:burst-akisi}). Sunucu tarafı, kareleri tek bir oturum (session)
altında biriktirir; 10.\ kare ulaştığında agregasyonu çalıştırır.

\begin{figure}[h!]
\centering
\begin{tikzpicture}[
    node distance=2mm, font=\scriptsize\ttfamily,
    every node/.style={draw=ink, minimum width=11mm, minimum height=7mm, align=center},
]
\node (n1) {f1\\\scriptsize sid=A,\,i=1/10};
\node (n2) [right=of n1] {f2};
\node (n3) [right=of n2] {f3};
\node (n4) [right=of n3] {f4};
\node (n5) [right=of n4] {f5};
\node (n6) [right=of n5] {f6};
\node (n7) [right=of n6] {f7};
\node (n8) [right=of n7] {f8};
\node (n9) [right=of n8] {f9};
\node (n10) [right=of n9, draw=signal, very thick] {f10\\\scriptsize final};
\end{tikzpicture}
\caption{ESP32-CAM tarafından yakalanan 10 karelik burst. Her kare aynı oturum kimliğini taşır;
yalnızca son kareye sunucu \emph{karar} ile yanıt verir, önceki karelerin yanıtı bir devam (continue) bildirimidir.}
\label{fig:burst-akisi}
\end{figure}

Bu yaklaşımın iki kazancı vardır:

\begin{enumerate}[label=(\arabic*),itemsep=2pt]
    \item \textbf{Robustluk:} 10 kareden bir kısmı kalite filtresine
          takılsa bile, en az 5 geçerli kare bulunduğunda doğruluk
          tek-kare durumuna kıyasla belirgin biçimde yükselir
          (literatür gözlemi: \cite{chen2018multiframe}).
    \item \textbf{Pasif canlılık:} 10 karelik gömme vektörlerinin
          karelere göre standart sapması, \emph{cross-frame embedding
          variance} olarak hesaplanır ve fotoğraf replay saldırılarını
          ek bir model yükü olmadan tespit etmek için kullanılır.
\end{enumerate}

\section{Sunucu Tarafı Kalite Filtresi ve Agregasyon}

Sunucu, her kare için aşağıdaki kalite kapısını uygular:

\begin{itemize}[itemsep=1pt]
    \item Yüz tespit skoru (\texttt{det\_score}) $\geq$ 0.7,
    \item Yüz bölgesi alanı (piksel) $\geq$ 80$\times$80,
    \item Yüz dönüklüğü (yaw) açısı $|$yaw$|$ $\leq$ 30$^{\circ}$,
    \item Laplace varyansı tabanlı bulanıklık skoru $\geq$ 10.
\end{itemize}

Bu kapıyı geçen karelerin embedding vektörleri biriktirilir. Son karede
en az \texttt{MIN\_QUALITY\_FRAMES} (varsayılan 5) geçerli kare yoksa,
sonuç \texttt{insufficient\_quality\_frames} olarak işaretlenir; ESP'ye
NOMATCH yerine kalite-yetersiz uyarısı döner. Yeterli kare varsa:

\begin{equation}
\vec{c} = \frac{1}{N}\sum_{i=1}^{N} \vec{e}_i, \quad
\vec{e}_i = \mathrm{ArcFace}(\mathrm{crop}(I_i)),
\end{equation}

şeklinde centroid hesaplanır ve veri tabanındaki her kayıt $\vec{d}_k$
ile kosinüs benzerliği ölçülür:

\begin{equation}
\mathrm{sim}_k = \frac{\vec{c} \cdot \vec{d}_k}{\|\vec{c}\|\, \|\vec{d}_k\|}.
\end{equation}

En yüksek benzerlik $\mathrm{sim}^\star = \max_k \mathrm{sim}_k$, eşik
$\tau$ (varsayılan 0.40) ile karşılaştırılır:
$\mathrm{sim}^\star \geq \tau \Rightarrow \text{MATCH}$. Pasif canlılık
skoru ise:

\begin{equation}
\mathrm{liveness} = \frac{1}{D}\sum_{j=1}^{D} \sigma_j, \quad
\sigma_j = \mathrm{std}\bigl(\{e_{i,j}\}_{i=1}^{N}\bigr),
\end{equation}

ile, embedding boyutlarına göre standart sapmaların ortalaması olarak
hesaplanır. Gerçek bir canlı yüzde bu skor sıfıra yakın ama anlamlı
biçimde sıfırın üstündedir; fotoğraf replay durumunda anlamlı biçimde
azalır.

\section{STM32 Olay Yöneticisi}
\label{sec:stm-olay-yonetimi}

STM32 NUCLEO-F767ZI üzerinde çalışan firmware, bir sonlu durum makinesi
(\emph{finite state machine}, FSM) etrafında düzenlenmiştir.
Modül HAL bağımlılığından arındırılmış, taşınabilir C99 olarak yazılmış
ve donanımla iletişimi bir \emph{callback vtable} üzerinden
gerçekleştirmektedir (\Cref{fig:fsm}). Bu sayede aynı kaynak dosyası
hem CubeIDE projesinde derlenebilmekte, hem de PC üzerinde sahte saat
ve sahte callback'lerle birim testten geçirilebilmektedir.

\begin{figure}[h!]
\centering
\begin{tikzpicture}[
    node distance=18mm and 18mm,
    state/.style={draw=ink, thick, rounded corners=2pt,
                  minimum width=20mm, minimum height=9mm,
                  align=center, font=\small\sffamily},
    accent/.style={draw=signal, very thick},
    arr/.style={-{Latex[length=1.6mm]}, thick, ink, font=\scriptsize\ttfamily},
]
\node[state] (idle) {IDLE};
\node[state, right=of idle] (scan) {SCANNING};
\node[state, above right=8mm and 14mm of scan] (match) {MATCH};
\node[state, below right=8mm and 14mm of scan] (nom) {NOMATCH};
\node[state, accent, below=20mm of scan] (panic) {PANIC};
\node[state, right=of nom] (neterr) {NETERR};

\draw[arr] (idle) -- node[above]{TARA} (scan);
\draw[arr] (scan) -- node[sloped, above]{MATCH} (match);
\draw[arr] (scan) -- node[sloped, below]{NOMATCH} (nom);
\draw[arr] (scan) -- node[sloped, above]{timeout / ERR} (neterr);
\draw[arr, dashed] (match.west) to[bend left=30] node[above]{hold 5\,s} (idle.north);
\draw[arr, dashed] (nom.west) to[bend right=30] node[below]{hold 2\,s} (idle.south);
\draw[arr, dashed] (neterr.north) to[bend left=15] node[right]{hold 3\,s} (idle.east);
\draw[arr, signal] (idle.south) to[bend right=25] node[left]{PANIC} (panic.west);
\draw[arr, signal] (scan.south) to[bend left=15] node[right]{PANIC} (panic.east);
\draw[arr, dashed] (panic.west) to[bend left=45] node[left]{hold 10\,s} (idle.south);
\end{tikzpicture}
\caption{STM32 olay yöneticisi durum geçiş diyagramı. PANIC olayı her durumda erişilebilir.}
\label{fig:fsm}
\end{figure}

Durum çıktıları aşağıdaki gibi tanımlanmıştır:

\begin{itemize}[itemsep=1pt]
    \item \textbf{IDLE:} yeşil LED sabit, 5\,s'de bir \texttt{HB\textbackslash n} kalp atışı,
    \item \textbf{SCANNING:} sarı LED sabit, UART1'e \texttt{CAPTURE\textbackslash n},
          UART2'ye \texttt{SCANNING\textbackslash n},
    \item \textbf{MATCH:} kırmızı LED + buzzer, HM-10'a \texttt{MATCH:<name>;<sim>\textbackslash n},
    \item \textbf{NOMATCH:} yeşil LED, HM-10'a \texttt{NOMATCH\textbackslash n},
    \item \textbf{PANIC:} kırmızı LED + buzzer, HM-10'a \texttt{PANIC\textbackslash n},
    \item \textbf{NETERR:} sarı LED yanıp sönüyor, HM-10'a \texttt{NETERR\textbackslash n}.
\end{itemize}

\section{Acil Çağrı Zinciri: BLE ve Android}

HM-10 modülü \emph{transparent UART$\leftrightarrow$BLE} köprüsü olarak
çalışır: STM32'nin UART2'ye yazdığı her bayt, modülün GATT karakteristiği
(\texttt{FFE0} servisi, \texttt{FFE1} karakteristiği) üzerinden eşli
telefona \emph{notification} olarak iletilir. Telefon tarafındaki TaxiGuard
uygulaması bir \emph{foreground service} olarak BLE central rolünde
durmakta, gelen ASCII satırları parse etmekte ve \texttt{MATCH:..} veya
\texttt{PANIC} olayı geldiğinde \texttt{Intent.ACTION\_CALL("tel:155")}
ile otomatik olarak Polis İmdat hattını aramaktadır.

\Cref{tab:mesaj-cercevesi} sistemde tanımlı mesaj çerçevelerini özetler.

\begin{table}[h!]
\centering
\caption{STM32 $\rightarrow$ HM-10 $\rightarrow$ Android mesaj çerçeveleri.}
\label{tab:mesaj-cercevesi}
\small
\begin{tabular}{l l l}
\toprule
\textbf{Mesaj} & \textbf{Tetik} & \textbf{Telefon eylemi}\\
\midrule
\texttt{SCANNING\textbackslash n} & TARA başlangıcı     & Banner: ``Tarama başladı''\\
\texttt{MATCH:<name>;<sim>\textbackslash n} & Pozitif eşleşme & \texttt{Intent.ACTION\_CALL("tel:155")}\\
\texttt{NOMATCH\textbackslash n}  & Burst negatif       & Banner: ``Eşleşme yok''\\
\texttt{PANIC\textbackslash n}    & PANİK butonu        & \texttt{Intent.ACTION\_CALL("tel:155")}\\
\texttt{NETERR\textbackslash n}   & Scan timeout/ERR    & Banner: ``Sunucuya erişilemedi''\\
\texttt{HB\textbackslash n}       & 5\,s aralıkla IDLE'da & Bağlantı sağlamlık göstergesi\\
\bottomrule
\end{tabular}
\end{table}

iOS yerine Android'in seçilme nedeni, iOS sandbox kısıtlamasının
\texttt{ACTION\_CALL} dengini desteklememesidir; iOS yalnızca
dialer'ı önceden doldurulmuş hâlde açabilen
\texttt{tel:} URL şemasına izin vermekte, otomatik bağlantı için
kullanıcı dokunması gerekmektedir.

\section{Ağ Yokken Davranış (Safe Mode)}
\label{sec:safe-mode}

Wi-Fi'nin yokluğunda veya sunucunun yanıt vermediği durumda sistem
\textbf{güvenli moda} geçer:

\begin{enumerate}[label=(\arabic*),itemsep=1pt]
    \item ESP32-CAM, HTTP POST sırasında 10~s zaman aşımına ulaşırsa
          UART1'e \texttt{ERR:no\_wifi} satırı gönderir.
    \item STM32 olay yöneticisi 15~s'lik scan timeout'una takılırsa
          (veya \texttt{ERR:..} satırı alırsa) NETERR durumuna geçer.
    \item NETERR durumunda sarı LED yanıp söner, buzzer susar ve HM-10'a
          \texttt{NETERR\textbackslash n} satırı gönderilir; telefon
          ``Sunucuya erişilemedi'' bildirimiyle sürücüyü uyarır.
    \item Bağlantı geri geldiğinde sistem normal akışa döner; ek bir
          yeniden bağlanma adımı gerektirmez.
\end{enumerate}

Sistem, ağ yokken sessizce başarısız olmaz; sürücüye görsel ve metinsel
geri bildirim ile durumu açıkça bildirir. Bu, hedef başarı ölçütleri
arasında tanımlanan ``ağ olmadığında güvenli mod'' gerekliliğini
karşılar.

\section{Öneri Raporuna ve Gelişme Raporuna Göre Değişiklikler}
\label{sec:degisiklikler}

Bu tezde, gelişme raporundaki tasarımdan iki önemli sapma vardır:

\begin{itemize}[itemsep=2pt]
    \item \textbf{Kamera platformu:} Plan A'da STM32'nin DCMI çevre
          birimine bağlı stand-alone bir OV5640 modülü öngörülmüştü.
          Sergi takvimi içinde DVP/parallel arabirimli OV5640 modülü
          tedarik edilemediği için \textbf{Plan B} olarak ESP32-CAM
          (AI-Thinker, OV3660 sensörü) kullanılmıştır. Bu değişikliğin
          \emph{tasarım üzerindeki etkisi:} kamera + Wi-Fi tek board'a
          taşınmış, STM32 yalnızca olay yöneticisi rolünde kalmıştır.
    \item \textbf{Gecikme hedefi:} Gelişme raporundaki ``$<$~200~ms''
          tek-kare ile çalışan bir API sorgusunu varsayıyordu. Çoklu kare
          burst tasarımına geçildiği için bu hedef, 10 kare için
          \emph{en fazla 7~s} olarak revize edilmiştir. Bu süre,
          yolcu binişi başına bir kez gerçekleşen bir işlem olduğu için
          kullanıcı deneyimi açısından uygundur.
\end{itemize}
